\section{Exceptional divisors between trace models}
\label{sec:trace-exceptional-divisors}

We return to the intrinsic trace data of $u$ along $Y$. Write $S_Z=S_Z(u;Y)$ on a smooth compact model $Z\to Y$. The finite-divisor inequality in \cref{thm:trace-divisorial-model-compatibility} now extends to the whole countable current. The difference is supported on finitely many exceptional divisors of the specified modification.

\begin{theorem}\label{thm:trace-exceptional-difference}
For a smooth compact modification $g:Z'\to Z$, there is a unique finite effective real divisor $E_g$ supported on its exceptional locus such that
\begin{equation}\label{eq:trace-exceptional-difference}
 S_{Z'}=g^*S_Z+[E_g].
\end{equation}
For $h:Z''\to Z'$, the divisors satisfy
\begin{equation}\label{eq:trace-exceptional-cocycle}
 E_{g\circ h}=E_h+h^*E_g.
\end{equation}
Their coefficients are additive and positively homogeneous in the ambient trace potential $u$.

The analytic-ideal approximants on the two models give finite effective exceptional divisors
\[
 [E_{m,g}]=D_{m,Z'}-g^*D_{m,Z}.
\]
For any Gauduchon metric $\gamma'$ on $Z'$, the measures
\[
 [E_{m,g}]\wedge(\gamma')^{r-1}
     \longrightarrow [E_g]\wedge(\gamma')^{r-1},\qquad r=\dim Y,
\]
converge in total variation.
\end{theorem}
\begin{proof}
Enumerate the positive divisorial components of $S_Z$. Every finite partial sum has pullback at most $S_{Z'}$ by \cref{thm:trace-divisorial-model-compatibility}. \Cref{cor:countable-divisorial-pullback} identifies the increasing limit of those pullbacks with $g^*S_Z$. Thus $S_{Z'}-g^*S_Z$ is a positive closed divisorial current. The two currents have equal coefficients along every strict transform of a divisor on $Z$, because these orders describe the same valuation. Every other prime divisor of $Z'$ is exceptional for $g$. The exceptional analytic set has only finitely many irreducible components on compact $Z'$. Its divisorial components therefore support a finite effective real divisor $E_g$. Coefficients of integration currents identify it uniquely.

Pullback of positive closed $(1,1)$-currents is functorial, by composition of local psh potentials. Using \eqref{eq:trace-exceptional-difference} for $h$ and $g$ gives
\[
 S_{Z''}=h^*S_{Z'}+[E_h]
        =(g\circ h)^*S_Z+[h^*E_g]+[E_h].
\]
Uniqueness gives \eqref{eq:trace-exceptional-cocycle}. Additivity and positive homogeneity follow from the same properties of every trace coefficient and linearity of pullback.

For a fixed $m$, the image ideal on $Z$ has a finite divisorial part. Its local generators are divisible by the corresponding Cartier equations. Extension to $Z'$ therefore gives
\[
 b_m(F)\geq\sum_{D\subset Z}
                    b_m(D)\operatorname{ord}_F(g^*D)
\]
for every prime divisor $F$ on $Z'$. The sum is finite. After division by $m$, this proves positivity of $D_{m,Z'}-g^*D_{m,Z}$. Its coefficients on strict transforms are zero, so it is the integration current of an effective exceptional real divisor $E_{m,g}$.

Choose also a Gauduchon metric $\gamma$ on $Z$. By \cref{thm:trace-divisorial-current-budget}, the currents $S_Z-D_{m,Z}$ and $S_{Z'}-D_{m,Z'}$ are positive and their trace masses tend to zero. Subtraction of the two identities for the exceptional currents gives
\[
 [E_{m,g}]-[E_g]
   =-(S_{Z'}-D_{m,Z'})+g^*(S_Z-D_{m,Z}).
\]
Both terms on the right, before the minus sign, are positive. The total variation of the signed trace measure is consequently at most
\[
 \int_{Z'}(S_{Z'}-D_{m,Z'})\wedge(\gamma')^{r-1}
 +C_g\int_Z(S_Z-D_{m,Z})\wedge\gamma^{r-1},
\]
by \eqref{eq:modification-linear-mass-bound}. It tends to zero. This proves the claimed convergence, without assuming monotonicity of the exceptional divisors as $m$ varies.
\end{proof}

The coefficient of $E_g$ at an exceptional prime $F$ is explicitly
\begin{equation}\label{eq:trace-exceptional-coefficient}
 \tau_{Y,u}(F)
  -\sum_{D\subset Z}\tau_{Y,u}(D)\operatorname{ord}_F(g^*D).
\end{equation}
The sum is countable and convergent. Indeed it is the coefficient of $g^*S_Z$ in the positive divisorial series of \cref{cor:countable-divisorial-pullback}; it is bounded above by $\tau_{Y,u}(F)$. This is a genuine whole-current formula, not a formal pullback of an unspecified infinite divisor.

\begin{corollary}\label{cor:trace-model-ordinary-comparison}
Let $v_Z$ and $v_{Z'}$ be the canonical ordinary-type representatives of $S_Z$ and $S_{Z'}$ from \cref{cor:divisorial-current-ordinary-type}. If $\rho_{E_g}$ is a global logarithmic potential of the finite real divisor $E_g$, then
\[
 v_{Z'}-(v_Z\circ g)-\rho_{E_g}
\]
is represented by a smooth function. In particular
\[
 v_{Z'}\preceq v_Z\circ g,
\]
and the two ordinary types are equal if and only if $E_g=0$.
\end{corollary}
\begin{proof}
Choose smooth Hermitian metrics on the line bundles of the prime components of $E_g$. The weighted sum of the logarithms of the squared canonical-section norms gives $\rho_{E_g}$ with
\[
 \ddc\rho_{E_g}=[E_g]-\alpha_{E_g},
\]
where $\alpha_{E_g}$ is smooth and closed. It is quasi-plurisubharmonic and bounded above on compact $Z'$. Writing $\ddc v_Z=S_Z-B_Z$ and $\ddc v_{Z'}=S_{Z'}-B_{Z'}$, the theorem gives
\[
 \ddc\bigl(v_{Z'}-v_Z\circ g-\rho_{E_g}\bigr)
       =g^*B_Z+\alpha_{E_g}-B_{Z'},
\]
which is smooth. The Hermitian-trace elliptic argument makes the $L^1$ difference smooth. Equality with the canonical qpsh representatives then holds everywhere. The bounded-above logarithmic potential gives the ordinary domination.

If $E_g=0$, the smooth difference gives ordinary equivalence. Conversely ordinary equivalence would give equal generic Lelong numbers on every prime divisor of $Z'$. The smooth residual terms show that their difference on a prime component of $E_g$ is precisely its positive coefficient, a contradiction. Hence ordinary equivalence forces $E_g=0$.
\end{proof}

When the original $\theta$ is closed, the residual Bott--Chern classes of \cref{prop:trace-bott-chern-bclass} satisfy the more precise relation
\[
 \mathbb A_{Z'}(u;Y,\theta)
       =g^*\mathbb A_Z(u;Y,\theta)-\{[E_g]\}_{\mathrm{BC}}.
\]
This follows by subtracting \eqref{eq:trace-exceptional-difference} from the pulled-back smooth background. Both residual classes are pseudoeffective, as already proved. The displayed formula describes their difference; it does not identify a nonlinear intersection product or volume.

\begin{example}\label{ex:exceptional-trace-with-exact-potential}
On $Y=\mathbb P^2$ put
\[
 u=a\log\frac{|Z_1|^2+|Z_2|^2}
                   {|Z_0|^2+|Z_1|^2+|Z_2|^2},\qquad a>0.
\]
Take the ambient manifold to be $Y$ itself. If $\sigma$ is the Fubini--Study form in our $\ddc$ normalization, then $u\in\PSH(Y,a\sigma)$. It has analytic singularities only at $p=[1:0:0]$, and has zero generic Lelong number on $Y$. Its trace is exactly the ordinary singularity data of $u$, by \cref{prop:trace-naive-comparison}. There are no divisorial poles on $Y$, so $S_Y=0$ and its canonical ordinary type is smooth. If $\pi:\widetilde Y\to Y$ is the blowup at $p$ with exceptional curve $F$, then
\[
 S_{\widetilde Y}=a[F],\qquad E_\pi=aF.
\]
Indeed on either standard blowup chart the numerator has one exceptional factor of squared modulus and a smooth positive remaining factor. The pole order is $a$ along $F$ and zero along every other prime divisor. Thus a nonzero exceptional difference does not obstruct the existence of an exact trace potential. It obstructs identifying that potential with the canonical type of the divisorial part on the original model.

The same distinction occurs on a manifold failing BMP. Take a compact Hermitian factor $(W,\omega_W)$ failing BMP, such as the Hopf threefold in \cref{thm:hopf-unbounded}, and pull this potential back to $Y\times W$. Its pole set is $\{p\}\times W$, of codimension two, and the modification $\widetilde Y\times W\to Y\times W$ has exceptional coefficient $a$ along $F\times W$. Analytic restriction again gives an exact trace. For a product Hermitian form $\Omega=\operatorname{pr}_Y^*\sigma+\operatorname{pr}_W^*\omega_W$, smooth mass witnesses $H_j$ on $W$ have pullback masses
\[
 \binom{k+2}{2}\left(\int_Y\sigma^2\right)
          \left(\int_W(\omega_W+\ddc H_j)^k\right)\longrightarrow+\infty,
 \qquad k=\dim W,
\]
by the smooth binomial expansion and the dimensions of the two factors. Hence the product also fails BMP. The potential is admitted by a sufficiently large Hermitian background; the BMP assertion does not impose an unproved positivity condition for a prescribed small background.
\end{example}

The exceptional divisors provide compatible information between every pair of models, and the analytic-ideal approximants recover them in linear mass. A general exact trace potential would still have to realize all higher-codimensional information simultaneously. The example shows why vanishing of every $E_g$ is not a necessary condition for that problem. It is precisely the condition for the canonical divisorial ordinary types themselves to pull back without an additional logarithmic term.
